\documentclass[a4paper,12pt]{article}
\usepackage[utf8]{inputenc}
\usepackage{graphicx}
\usepackage{booktabs}
\usepackage{amsmath}
\usepackage{hyperref}
\usepackage{listings}
\usepackage{geometry}
\geometry{left=2.5cm, right=2.5cm, top=3cm, bottom=3cm}

\title{Trabajo de Econometría}

\date{}


\begin{document}

\maketitle

\section{Introducción}

El objetivo de este trabajo es aplicar herramientas econométricas utilizando el software \textbf{R} y las bases de datos del paquete \texttt{wooldridge}. Se busca desarrollar habilidades en el análisis de datos, estimación de modelos de regresión y evaluación de supuestos econométricos.

El trabajo se realizará en grupos de 5 o 6 integrantes. Cada grupo deberá elegir una base de datos del paquete \texttt{wooldridge} de la siguiente lista:

\begin{itemize}
    \item \textbf{attend}: Asistencia a clases y su impacto en el rendimiento académico.
    \item \textbf{bwght2}: Factores que afectan el peso al nacer de los bebés.
    \item \textbf{ceosal1}: Relación entre los salarios de CEOs y el desempeño de las empresas.
    \item \textbf{beauty}: Impacto de la apariencia física en los salarios.
    \item \textbf{hprice2}: Factores que determinan el precio de las viviendas.
    \item \textbf{mlb1}: Salarios de jugadores de béisbol de la MLB y su desempeño.
    \item \textbf{vote1}: Resultados electorales y factores políticos en distritos de EE.UU.
    \item \textbf{injury}: Datos sobre lesiones laborales y pagos de compensación.
    \item \textbf{GPA2}: Factores que influyen en el rendimiento académico universitario.
    \item \textbf{nbasal}: Salarios de jugadores de la NBA y su desempeño.
    \item \textbf{htv}: Educación y decisiones en el mercado laboral.
\end{itemize}

\section{Estructura y Evaluación}

El trabajo consta de tres etapas, con las siguientes ponderaciones:

\begin{itemize}
    \item \textbf{Avance preliminar (10\% Talleres de ayudantía)}
    
    Se debe entregar un documento inicial en \texttt{Google Colab} con:
    \begin{itemize}
        \item Descripción de la base de datos seleccionada, incluyendo su origen y principales variables.
        \item Medidas de tendencia central y dispersión de las variables clave.
        \item Gráficos descriptivos de las variables más relevantes.
        \item Planteamiento de las hipótesis de investigación.
        \item Justificación de la elección de la variable dependiente y explicativas.
        \item Especificación del modelo de regresión poblacional propuesto.
    \end{itemize}
    
    \item \textbf{Informe final (12,5\%)}

    Se debe entregar un informe en formato PDF, acompañado de un \texttt{Google Colab} con los códigos en \texttt{R} que permitan replicar los resultados (incluyendo explicaciones breves de los comandos utilizados).
    
    El informe debe incluir:
    
    \begin{itemize}
        \item \textbf{Estimación del modelo}: 
        \begin{itemize}
            \item Análisis de los coeficientes: signo, magnitud e interpretación económica.
            \item Evaluación de la significancia estadística (\textit{p-values}).
            \item Medidas de bondad de ajuste (\textit{R-cuadrado}).
        \end{itemize}

      

        \item \textbf{Verificación de los supuestos del Modelo de Regresión Lineal Múltiple (MCRL)}:
        \begin{itemize}
            \item \textbf{Exogeneidad de los regresores}: El término de error debe cumplir:
            \[
            E(u_i | X_{i1}, X_{i2}, \dots, X_{ik}) = 0.
            \]
            Esto garantiza que los regresores no están correlacionados con el error.

            \item \textbf{Homoscedasticidad} (Breusch-Pagan, White)
            
            \item \textbf{Normalidad de los errores}
                      \item \textbf{Multicolinealidad} (Factor de Inflación de la Varianza - VIF):
           
        \end{itemize}

        Si se detectan problemas en los supuestos, se deben proponer y aplicar correcciones.
          \item \textbf{Testear hipótesis de la investigación }:

        \item \textbf{Conclusiones y recomendaciones}:
        \begin{itemize}
            \item Resumen de los principales hallazgos del análisis.
            \item Evaluación de si las hipótesis fueron confirmadas o rechazadas.
            \item Limitaciones del modelo y propuestas de mejoras.
        \end{itemize}
    \end{itemize}

    \item \textbf{Presentación oral (12,5\%)}

    Cada grupo dispondrá de \textbf{15 minutos} para exponer sus resultados. Un estudiante será seleccionado aleatoriamente para presentar, y luego todo el grupo podrá participar en la ronda de preguntas.

    Se evaluará:
    \begin{itemize}
        \item Organización y claridad de la presentación.
        \item Uso adecuado del lenguaje técnico.
        \item Capacidad de responder preguntas y defender los resultados obtenidos.
    \end{itemize}
\end{itemize}

\end{document}
